\section{实验}
\subsection{数据集与评测协议}
本文采用 RSKT-Seg 的八数据集评测集合~\cite{liExploringEfficientOpenVocabulary2026}，顺序固定为 DLRSD、iSAID、LoveDA、Potsdam、Vaihingen、UAVid、UDD5 和 VDD。该集合覆盖地表覆盖、城市场景、密集航空目标及无人机斜视图像。集合相同不意味着划分、类别映射和处理方式已经一致；本地 RSKT-Seg 正式版将详细说明指向未随附件提供的附录，因而本稿暂不声称完全复用其评测协议。表~\ref{tab:datasets_protocol} 列出本地配置及待核实项。

主设置为 iSAID 单源监督适配，iSAID 同时用于源域保持诊断，其余七个数据集构成预定目标集合。DLRSD 单源仅作为可选扩展，若实施则另列 panel，不与 iSAID 训练结果混合。当前开发入口在 LoveDA val 上以含背景的七类 mIoU 选择 checkpoint；因此该划分参与了模型选择，不能作为未参与选择的独立目标评测。\TODO{正式运行前固定源域留出验证或预先固定训练步数等不使用目标标签的选择方案，并核实原始大图或视频序列之间的独立性；本轮未实施该调整。若保留 LoveDA val 选择，须将其从独立目标均值中排除。}

主指标为平均交并比 mIoU（\%），补充平均类别准确率 mAcc 和逐类 IoU。LoveDA 仅使用包含 background 的七类口径；逐类保存 background、building、road、water、barren、forest 和 agricultural 的 IoU。其他数据集按核实后的基准定义处理，不能统一删除背景。零并集类别政策、有效类别数与标签映射均需随结果保存；未知数据不记为零。源域与目标域的类别语义重叠、基础模型预训练数据是否覆盖目标图像仍需披露，跨数据集迁移不能自动称为严格未见类别评测。

Avg-8 仅对八列齐备且协议已解释的结果等权计算，它包含源域，不能称跨域平均。Target Avg 的预定共同集合为除 iSAID 外的七个数据集；在当前选择设置下，它包含开发域 LoveDA。若报告 Independent Target Avg，其集合固定为 DLRSD、Potsdam、Vaihingen、UAVid、UDD5 和 VDD。缺少任一必需单元格或协议未锁定时均显示 --，不按每行可用项改变分母，不复制原论文的跨数据集平均。文献值保留原报告精度，自测值及派生均值预定保留两位小数；变化量用百分点表达。
\begin{table*}[t]
\centering
\caption{八数据集协议登记。类别数和背景规则来自本地配置，并不表示已与 RSKT-Seg 重处理版本一致。所有本地 ignore 值为 255；样本清单及原始场景独立性待核实。}
\label{tab:datasets_protocol}
\small
\setlength{\tabcolsep}{3pt}
\begin{tabular}{lL{0.11\textwidth}L{0.17\textwidth}L{0.19\textwidth}L{0.27\textwidth}}
\toprule
数据集 & 本地类别数 & 本地评测划分 & 背景规则／角色 & 原始数据来源与基准状态\\
\midrule
DLRSD & 17 & val & 无背景类／目标 & \cite{PerformanceEvaluationSingleLabel}；基准类别与划分 TODO\\
iSAID & 15 & val & 原始零标签忽略／源域 & \cite{zamirISAIDLargescaleDataset2019}；训练为 train，基准版本 TODO\\
LoveDA & 7 & val & 含背景／开发选择域 & \cite{wangLoveDARemoteSensing2022}；主指标固定含背景七类，独立正式划分 TODO\\
Potsdam & 6 & val & clutter 保留评测／目标 & 原始官方资源条目 TODO；RSKT-Seg 具体重处理 TODO\\
Vaihingen & 6 & val & clutter 保留评测／目标 & 原始官方资源条目 TODO；RSKT-Seg 具体重处理 TODO\\
UAVid & 8 & val & 含背景／目标 & \cite{lyuUAVidSemanticSegmentation2020}；需核对车辆合并及序列独立性\\
UDD5 & 5 & val & 含背景／目标 & \cite{chenLargeScaleStructureMotion2018}；UDD5 版本与标签映射 TODO\\
VDD & 7 & val & 含背景／目标 & \cite{caiVDDVariedDrone2025} 原始全文缺失；基准版本 TODO\\
\bottomrule
\end{tabular}
\par\vspace{3pt}
\begin{minipage}{\textwidth}\footnotesize
八数据集集合来源为 RSKT-Seg~\cite{liExploringEfficientOpenVocabulary2026}，原始数据引用用于数据介绍，不能代替基准划分证据。本地 iSAID 排除背景，而 SegEarth-OV 本地版本使用含背景定义；UAVid 与 UDD5 在其他文献中存在类别合并或删背景设置，不能直接认作同协议。LoveDA 为待定的独立性例外，不标成独立目标评测。
\end{minipage}
\end{table*}

% Reference table; not included in the main manuscript.
\begin{table}[!t]
\centering
\caption{Image and label preprocessing.}
\label{tab:dataset_processing}
\small
\begin{tabular}{lL{0.60\columnwidth}}
\toprule
Item & Setting\\
\midrule
Image size & $1008\times1008$ for SAM 3; $504\times504$ for RemoteCLIP\\
Label handling & Dataset-specific background and ignore mappings\\
Evaluation & Full validation split; no test-time augmentation\\
\bottomrule
\end{tabular}
\end{table}



\subsection{实现细节与比较设置}
当前代码采用冻结 SAM 3 编码器、Pixel Decoder 与语义头的适配路径，训练 RemoteCLIP 引导特征聚合及 Semantic--Detail Fusion（SDF，语义—细节融合）模块，并部分微调 RemoteCLIP 图文分支的注意力 Q/V 和位置参数。冻结解码器在学生路径中仍允许输入梯度通过，不能将整个系统称为完全冻结。表~\ref{tab:implementation} 的数值来自当前合并配置，只表示待锁定的实现起点，不能作为已执行实验的日志。

文献成绩与本文实测分开登记。本文数值尚未确认；文献数值逐项记录实际本地 PDF 的版本、哈希、页码、表号和行列，方法原论文与数字来源分别引用。资源分组依据原文训练数据和标签使用情况，而非是否冻结参数。所有自测比较将锁定训练数据、模型选择、输入、提示、损失基准、推理和硬件；文献训练预算未统一，不能写成全部方法在同一预算下复现。
% Reference table; implementation settings are described in the main text.
\begin{table}[!t]
\centering
\caption{Implementation settings of ReCA-SAM3.}
\label{tab:implementation}
\small
\begin{tabular}{lL{0.62\columnwidth}}
\toprule
Component & Setting\\
\midrule
Backbones & Frozen SAM 3; RemoteCLIP ViT-L/14\\
Aggregation & Four blocks; eight heads; window/shift 12/6\\
Fusion & Two 256-channel branches; four orientations\\
Optimizer & AdamW; base learning rate $10^{-4}$\\
Training & Batch size 2; 20,000 steps; 1,000 warmup steps\\
Loss weights & Segmentation 1; distillation 0.5; suppression 0.01\\
Inference & Prompt chunk size 4; no test-time augmentation\\
\bottomrule
\end{tabular}
\end{table}



\subsection{与现有方法的比较}
表~\ref{tab:main_rskt8} 将本文待测行与文献参考区域分开，资源与来源见表~\ref{tab:resource_settings}。候选包括 SAM 3 Semantic-only、SegEarth-OV、完整 SegEarth-OV3、ProC-SAM3、CAT-Seg、RSKT-Seg、Pi-Seg、GSNet，以及本文的简单 adapter 和 ReCA-SAM3。SAM 3 Semantic-only 的已有局部数值来自 SegEarth-OV3 对仅语义头变体的实验，不能替代完整 SegEarth-OV3；GSNet 的 iSAID 行是 RSKT-Seg 对该方法的重评设置，不能把 GSNet 原论文的 LandDiscover50K 训练改称 iSAID 训练。

本地附件还揭示了版本与复现差异。SegEarth-OV 和 ProC-SAM3 的实际附件分别为 arXiv v2，而文献库条目分别登记 CVPR 和 GRSL 发表信息；本稿数值仅归属于实际读取版本。SegEarth-OV3 原论文与 ProC-SAM3 标星复现行分别存档，未选取较高值或拼接为一行。LoveDA 类别数和背景口径未获该来源明确确认的单元格不进入七类列。当前没有已确认与本项目完全匹配的文献行，因此参考区不排名、不加粗；其数值不证明本文已经取得优势。Pi-Seg 的多源设置、GSNet 原始大规模训练等资源差异另作披露，不能混入 iSAID 单源行。PCA-Seg 保留在相关工作；SkySense-O 等未完成协议对齐的可选方法不强行补入主表。

\TODO{在协议与自测结果齐备后，比较源域保持、共同目标集合上的迁移及失败情形，结合训练资源解释差异。当前不写最优、显著提升或已证明泛化。}
% Generated from paper/data; do not hand-edit numeric cells.
\begin{table*}[t]
\centering
\caption{八数据集 mIoU（\%）框架与文献参考结果。iSAID 为本文源域；其余列为目标数据集。-- 表示未完成、未报告或口径未核实，所有数值均不排名。}
\label{tab:main_rskt8}
\small
\setlength{\tabcolsep}{3pt}
\begin{tabular}{L{0.245\textwidth}cccccccccc}
\toprule
方法 & DLRSD & iSAID & LoveDA & Potsdam & Vaihingen & UAVid & UDD5 & VDD & Avg-8 & Target Avg\\
\midrule

\multicolumn{11}{l}{本文计划结果（尚未完成）}\\
ReCA-SAM3$^{N}$ & -- & -- & -- & -- & -- & -- & -- & -- & -- & --\\
简单 SAM 3 adapter$^{N}$ & -- & -- & -- & -- & -- & -- & -- & -- & -- & --\\
\midrule
\multicolumn{11}{l}{文献参考结果（资源或协议尚未匹配，不参与排名）}\\
SAM 3 Semantic-only~\cite{carionSAM3Segment2026}$^{S}$ & -- & -- & -- & -- & -- & 47.1 & -- & -- & -- & --\\
SegEarth-OV~\cite{liSegEarthOVTrainingFreeOpenVocabulary2025}$^{E}$ & -- & 21.7 & 36.9 & 47.1 & 29.1 & 42.5 & 50.6 & 45.3 & -- & --\\
SegEarth-OV3~\cite{liSegEarthOV3ExploringSAM2026}$^{S}$ & -- & 27.6 & -- & 57.8 & 60.8 & 54.7 & 71.7 & 64.5 & -- & --\\
ProC-SAM3~\cite{songPromptCalibratedSAM32026}$^{P}$ & -- & 29.5 & -- & 58.8 & 64.5 & 60.0 & 73.4 & 67.5 & -- & --\\
CAT-Seg / iSAID~\cite{choCATSegCostAggregation2024}$^{R}$ & 28.80 & 94.77 & -- & 23.90 & 21.74 & 16.10 & 24.32 & 30.16 & -- & --\\
RSKT-Seg / iSAID~\cite{liExploringEfficientOpenVocabulary2026}$^{R}$ & 31.57 & 93.96 & -- & 28.57 & 25.55 & 17.36 & 34.01 & 33.40 & -- & --\\
Pi-Seg / iSAID~\cite{liPerturbationRegularizedOpenVocabularyRemote2026}$^{I}$ & 33.67 & -- & -- & 28.14 & 23.83 & 12.34 & 37.24 & 38.24 & -- & --\\
GSNet / iSAID re-eval.~\cite{yeOpenVocabularyRemoteSensing2025}$^{R}$ & 31.50 & 93.11 & -- & 28.50 & 25.10 & 15.98 & 32.24 & 32.07 & -- & --\\
\bottomrule
\end{tabular}
\par\vspace{3pt}
\begin{minipage}{\textwidth}
\footnotesize N：本文待测；S、E、P、R、I 为表~\ref{tab:resource_settings} 的资源/来源标记。R 的数值来自 RSKT-Seg~\cite{liExploringEfficientOpenVocabulary2026} Table~1 的 iSAID/ViT-L 区；S 的数值来自 SegEarth-OV3~\cite{liSegEarthOV3ExploringSAM2026} Table~1 或 Table~4；E 为 SegEarth-OV 本地 arXiv v2 Table~1；P 为 ProC-SAM3 本地 arXiv v2 Table~I；I 为 Pi-Seg arXiv v3 Table~1 的 iSAID/ViT-L 区。LoveDA 仅展示已确认的含背景七类成绩。平均集合与协议未锁定，均值保留 --，不复制源论文的平均。
\end{minipage}
\end{table*}

% Reference table; detailed source auditing is maintained outside the manuscript.
\begin{table*}[!t]
\centering
\caption{Training resources and numerical sources of the comparison methods.}
\label{tab:resource_settings}
\small
\setlength{\tabcolsep}{4pt}
\begin{tabular}{L{0.19\textwidth}L{0.34\textwidth}L{0.35\textwidth}}
\toprule
Method & Training resources & Numerical source\\
\midrule
SAM 3 Semantic-only / SegEarth-OV3 & SAM 3; no downstream parameter adaptation & SegEarth-OV3, arXiv v2, Tables 1 and 4\\
SegEarth-OV & CLIP ViT-B/16; unlabeled Million-AID images for SimFeatUp & SegEarth-OV, CVPR 2025, Table 1\\
ProC-SAM3 & SAM 3; offline prompt generation & ProC-SAM3, arXiv v2, Table I\\
CAT-Seg / GSNet / RSKT-Seg & iSAID single-source adaptation; ViT-L panel & RSKT-Seg, AAAI version, Table 1\\
Pi-Seg & iSAID single-source adaptation; ViT-L panel & Pi-Seg, arXiv v3, Table 1\\
ReCA-SAM3 / adapter & DLRSD single-source adaptation & Native evaluation\\
\bottomrule
\end{tabular}
\end{table*}



\subsection{核心模块消融}
核心矩阵包含十个训练设置 T00--T09，表~\ref{tab:ablation_core} 与表~\ref{tab:ablation_losses} 共用 T00 Full，不为每张表重复训练。结构及损失变体从相同预训练起点重新训练，不能在完整权重上临时关闭模块当作训练消融。优先在 LoveDA 七类、Potsdam、UAVid 报告，Full、adapter 和 T02 优先覆盖全部八数据集；iSAID 另用于源域诊断。先用一个预先固定的 seed 完成矩阵，Full、adapter、T02 优先补三个预先指定 seed；当前正式种子清单为 TODO，代码默认 seed 不代表已完成运行，单 seed 不显示虚构标准差。

T00 保留完整当前设计。T01 使用所有提示共享的残差 adapter：原始 256 通道提示条件特征依次经过 $1\times1$ 卷积（256 到 256）、GroupNorm（GN，组归一化）、GELU、$3\times3$ 卷积（256 到 256）、GN，再与原输入残差相加，送入相同冻结解码器；不使用 RemoteCLIP、特征聚合和 SDF。其成本单列，不声称与 Full 参数完全相同。

T02 在所有类间块的两次注意力中屏蔽展开提示轴的全部非对角项，仅允许每个提示关注自身，保留投影和残差。该设置同时阻断同类别别名之间的信息交换，不能保留别名交互却称完全无跨提示交互。T03 bypass 两条流的所有类内块，保持类间流程。T04 将图像、文本和中间层指导全部替换为同 ViT-L/14 架构的通用 CLIP 权重，保持提取和部分微调策略。

T05 将初始语义状态设为同形状零张量，保留 CLIP 中间层指导和双流结构，因此称“无初始模板分数先验”，不能称“无 RemoteCLIP”。T06 将上采样聚合特征、原始提示条件特征和 FPN（特征金字塔网络）三路各 256 通道拼接，依次经 $1\times1$ 卷积（768 到 256）、GN、GELU 和 $3\times3$ 卷积（256 到 256）得到相同解码器输入；这是待实现对照，尚未声称代码支持。T07 保留 SDF 两支路的 256 通道、投影、归一化及残差，仅执行 $0^\circ$ 卷积，不与旧 128 通道版本比较。

T08 仅将参考蒸馏权重设零，T09 仅将忽略区域抑制权重设零，其他监督保持不变。它们回答监督约束的贡献，不扩展为大量损失权重扫描。所有开关定义和执行状态记录于实验计划；未实现变体不生成伪配置，所有成绩均待测。
% Generated from paper/data; do not hand-edit numeric cells.
\begin{table*}[t]
\centering
\caption{核心结构消融计划。所有成绩待测；每行定义见正文，结构变体需重新训练。}
\label{tab:ablation_core}
\small
\setlength{\tabcolsep}{3pt}
\begin{tabular}{L{0.30\textwidth}ccc}
\toprule
变体 & LoveDA & Potsdam & UAVid\\
\midrule
T00 完整 ReCA-SAM3 & -- & -- & --\\
T01 简单 SAM 3 adapter & -- & -- & --\\
T02 阻断跨类别交换 & -- & -- & --\\
T03 无类内聚合 & -- & -- & --\\
T04 通用 CLIP & -- & -- & --\\
T05 无初始模板分数先验 & -- & -- & --\\
T06 简单拼接融合 & -- & -- & --\\
T07 SDF 单方向 & -- & -- & --\\
\bottomrule
\end{tabular}
\par\vspace{3pt}
\begin{minipage}{\textwidth}
\footnotesize mIoU 单位为百分比；当前 LoveDA 为开发选择域，独立正式划分仍待确认；种子计划未锁定，不生成标准差。
\end{minipage}
\end{table*}

% Generated from paper/data; do not hand-edit numeric cells.
\begin{table*}[t]
\centering
\caption{监督目标消融计划。Full 为同一共享 T00；其余每行仅关闭一个损失项。}
\label{tab:ablation_losses}
\small
\setlength{\tabcolsep}{3pt}
\begin{tabular}{L{0.30\textwidth}ccccc}
\toprule
变体 & 蒸馏 & 抑制 & LoveDA & Potsdam & UAVid\\
\midrule
T00 完整 ReCA-SAM3 & 开 & 开 & -- & -- & --\\
T08 无参考蒸馏 & 关 & 开 & -- & -- & --\\
T09 无忽略区域抑制 & 开 & 关 & -- & -- & --\\
\bottomrule
\end{tabular}
\par\vspace{3pt}
\begin{minipage}{\textwidth}
\footnotesize mIoU 单位为百分比；当前 LoveDA 为开发选择域，独立正式划分仍待确认；种子计划未锁定，不生成标准差。
\end{minipage}
\end{table*}


\subsection{推理后处理分析}
推理双门采用同一 checkpoint 的四组评测，见表~\ref{tab:postprocess_gates}：P00 两门关闭，P10 仅开启原始 SAM 3 reference 门，P01 仅开启学生 prediction 门，P11 两门均开启。关闭使用实际代码的 None 语义；开启阈值、背景回填阈值、测试时增强（TTA）、权重哈希与实际推理时间单独登记。reference 门需要额外原始 SAM 3 解码，其成本计入端到端时间；四组不增加训练次数。

主训练消融采用同一预先锁定的推理方案；阈值若需选择，只在开发划分确定，不针对每个正式目标集取最优值。\TODO{为 Full 与 T02 补充未过滤输出，检查结构差异是否被后处理掩盖；尚无该证据，不预写收益结论。本轮不增加阈值扫描曲线。}
% Generated from paper/data; do not hand-edit numeric cells.
\begin{table}[!tb]
\centering
\caption{Effect of inference gates (mIoU, \%).}
\label{tab:postprocess_gates}
\small
\setlength{\tabcolsep}{2pt}
\begin{tabular}{lcccccc}
\toprule
ID & \shortstack{Ref.\\gate} & \shortstack{Pred.\\gate} & LoveDA & Potsdam & UAVid & \shortstack{Time\\(ms)}\\
\midrule
P00 & Off & Off & -- & -- & -- & --\\
P10 & On & Off & -- & -- & -- & --\\
P01 & Off & On & -- & -- & -- & --\\
P11 & On & On & -- & -- & -- & --\\
\bottomrule
\end{tabular}
\par\vspace{3pt}
\begin{minipage}{\columnwidth}
\footnotesize All settings share one checkpoint. Ref. and Pred. denote reference and prediction gates. Latency includes reference decoding when enabled. -- denotes unavailable results.
\end{minipage}
\end{table}


\subsection{定性比较与计算开销}
图~\ref{fig:qualitative_comparison} 规划在相同样本、色表和局部放大区域展示图像、真值及实际预测，同时包含成功与失败案例。只有确实获得本地输出的方法才展示像素预测，不能根据文献表格生成他人的预测图，也不将原论文示意图冒充本地结果。
\begin{figure}[t]
\centering
\fbox{\begin{minipage}[c][2.4cm][c]{0.92\columnwidth}
\centering
Qualitative Comparison Placeholder\\[0.4em]
Image / Ground Truth / SAM 3\\
Adapter / ReCA-SAM3\\[0.4em]
Matched Samples, Insets, and Failure Cases
\end{minipage}}
\caption{定性比较规划占位。正式图内使用英文标注；样本、局部放大和失败案例在获得真实预测后固定，当前无预测结果。}
\label{fig:qualitative_comparison}
\end{figure}
表~\ref{tab:efficiency} 分列总参数、可训练参数、端到端时间与峰值显存。固定 GPU、精度、输入、batch、提示数量、chunk、缓存状态和双门，并记录原始 SAM 3 reference 解码的额外成本；不能用不同文献硬件上的耗时作公平排名。本轮尚无确认的成本测量。
% Generated from paper/data; do not hand-edit numeric cells.
\begin{table}[t]
\centering
\caption{计算开销计划。总参数与可训练参数分列，时间和峰值显存采用统一测量设置。}
\label{tab:efficiency}
\small
\setlength{\tabcolsep}{3pt}
\begin{tabular}{L{0.31\columnwidth}cccc}
\toprule
方法 & \shortstack{总参数\\(M)} & \shortstack{可训练\\(M)} & \shortstack{时间\\(ms)} & \shortstack{显存\\(GB)}\\
\midrule
SAM 3 Semantic-only & -- & -- & -- & --\\
简单 adapter & -- & -- & -- & --\\
T00 Full & -- & -- & -- & --\\
T07 SDF 单方向 & -- & -- & -- & --\\
\bottomrule
\end{tabular}
\par\vspace{3pt}
\begin{minipage}{\columnwidth}
\footnotesize 固定 GPU、精度、输入、batch、提示数、chunk、缓存状态与双门。Reference 解码成本计入端到端时间；不将异硬件文献耗时用于公平排名。
\end{minipage}
\end{table}


\subsection{进一步分析}
优先诊断 Full 与 T02 对上下文提示集合的响应：固定最终输出类别，额外上下文仅参与聚合时单列诊断协议，避免改变 argmax 候选数造成混淆；该模式尚待实现。小目标分析拟在 iSAID 按实例面积诊断语义召回，它不是实例 AP；LoveDA/Potsdam 的混淆矩阵需匹配类别本体，忽略区域不当作完整负例。旋转一致性、边界和严格 adaptation-unseen 类别评测仅为可选后续，需显式类别划分和监督隔离后再作结论。
\FloatBarrier
